% Generated draft: requires mathematical and source review.
\documentclass[11pt]{article}
\usepackage[margin=25mm]{geometry}
\usepackage{fontspec}
\setmainfont{Latin Modern Roman}
\usepackage{amsmath,amssymb,mathtools,unicode-math}
\setmathfont{Latin Modern Math}
\usepackage{enumitem,needspace,xurl,hyperref}
\setlength{\parindent}{0pt}
\setlength{\parskip}{5pt}
\setlength{\emergencystretch}{3em}
\newcommand{\N}{\mathbb{N}}
\newcommand{\Z}{\mathbb{Z}}
\newcommand{\Q}{\mathbb{Q}}
\newcommand{\R}{\mathbb{R}}
\newcommand{\C}{\mathbb{C}}
\newcommand{\F}{\mathbb{F}}
\newcommand{\A}{\mathbb{A}}
\newcommand{\PP}{\mathbb{P}}
\newcommand{\E}{\mathbb{E}}
\newcommand{\eps}{\varepsilon}
\newcommand{\dd}{\,\mathrm{d}}
\begin{document}
\section*{UnsolvedMath ID: 700432}
\textbf{Source problem code:} catalogue entry
\subsection{Definitions and mathematical statement}
A variety of groups \(\mathfrak V\) is the class of all groups satisfying a specified set of group identities. Let \(W_{\mathfrak V}\) be the set of group words whose associated identities \(w=1\) hold in every group in \(\mathfrak V\). For any group \(G\), define its verbal subgroup by
\[
\mathfrak V(G)=\left\langle w(g_1,\ldots,g_n)\;\middle|\; n\geq 1,\ w\in W_{\mathfrak V}\text{ a word in }n\text{ variables},\ g_1,\ldots,g_n\in G\right\rangle.
\]
This subgroup is normal in \(G\), and \(G/\mathfrak V(G)\) is the largest quotient of \(G\) belonging to \(\mathfrak V\). The supplied source does not define \(\mathfrak V(G)\); the standard verbal-subgroup interpretation is used here.

A group \(H\) is residually finite if, for every \(h\in H\setminus\{1\}\), there exist a finite group \(F\) and a homomorphism \(\varphi:H\to F\) such that \(\varphi(h)\neq 1\).

\textbf{Problem.} Characterize all varieties of groups \(\mathfrak V\) such that
\[
\text{for every group }G,\qquad G\text{ residually finite}\ \Longrightarrow\ G/\mathfrak V(G)\text{ residually finite}.
\]
The quantifier ranges over arbitrary groups; the supplied statement imposes no finite-generation or cardinality restriction.
\subsection{Short English statement}
Determine which varieties of groups have the property that the largest quotient belonging to the variety of every residually finite group is again residually finite. Here \(\mathfrak V(G)\) is interpreted as the verbal subgroup associated with the variety.
\subsection{Sources}
\begin{itemize}
\item Local UnsolvedMath registry, supplied complete record with permanent ID \(700432\), exact source code \texttt{catalogue entry}, title \emph{Kourovka Notebook Problem 7.35}. The registry labels the problem open; this is supplied source metadata, not an independently verified assessment. Supplied upstream revision: \texttt{d371ee917ac9647fbb06cd6fc529deec6f3dda82}.
\item \emph{The Kourovka Notebook (Unsolved Problems in Group Theory)}, Problem 7.35, proposed by O. V. Mel'nikov in the 7th issue (1980). The supplied background references the 21st issue edition (2026), arXiv:1401.0300v48. The supplied collection description names E. I. Khukhro and V. D. Mazurov as editors.
\end{itemize}
\end{document}
